\documentclass[14pt]{extarticle}
\usepackage{fontspec}
\usepackage{polyglossia}
\usepackage{amsmath}
\usepackage{amssymb}
\usepackage{xcolor}
\usepackage{fancyhdr}
\usepackage{graphicx}
\usepackage{listings}
\usepackage[most]{tcolorbox}
% \usepackage{setspace}
\usepackage{pifont}
\usepackage{enumitem}
\usepackage{titlesec}
\usepackage[bottom]{footmisc}
\usepackage{titling}
\usepackage{minted}
\usepackage{etoolbox}
\usepackage{array}
\usepackage{longtable} % page-breakable tables (must load before polyglossia/bidi RTL setup)
\usepackage{extsizes}
\usepackage{float}
\usepackage{booktabs}
\usepackage{hyperref}
\usepackage[onehalfspacing]{setspace}
\newif\ifwhiteblack
\ifdefined\setwhiteblack
  \setwhiteblack
\fi

\ifwhiteblack
  \usepackage[a4paper, landscape, top=1.25cm, bottom=1.25cm, left=1.25cm, right=1.25cm, headheight=1.5cm, headsep=0.5cm]{geometry}
\else
  \usepackage[a4paper, top=2.7cm, bottom=2.7cm, left=2.5cm, right=2.5cm, headheight=1.5cm, headsep=0.5cm]{geometry}
\fi
\usepackage{tikz}
\usepackage{datetime2}

\usetikzlibrary{automata, positioning, arrows.meta, arrows}
\providecommand{\LTcaptype}{table}

\newfontfamily\emoji{Segoe UI Emoji}

\pagestyle{fancy}

\setmainlanguage[numerals=western]{arabic}
\setotherlanguage{english}
\newfontfamily\arabicfont[Script=Arabic]{Calibri}
\newfontfamily\arabicfonttt[Script=Arabic]{Courier New}



\ifwhiteblack
\lstset{
  language=[Sharp]C,
  numbers=left,
  stepnumber=1,
  numbersep=8pt,
  frame=single,
  basicstyle=\ttfamily\small,
  keywordstyle=\bfseries,
  stringstyle=\ttfamily,
  commentstyle=\itshape
}
\else
\lstset{
  language=[Sharp]C,
  numbers=left,
  stepnumber=1,
  numbersep=8pt,
  frame=single,
  basicstyle=\ttfamily\small,
  keywordstyle=\color{blue},
  stringstyle=\color{red},
  commentstyle=\color{green!50!black}
}
\fi

\newif\ifdetailed
\ifdefined\setdetailed
  \setdetailed
\fi

\newif\ifwithsols
\ifdefined\setwithsols
  \setwithsols
\fi

\ifwhiteblack
  \AtBeginDocument{
    \LARGE
    \usemintedstyle{bw}
  }
\fi

% Default minted look: compact code that wraps long lines.
% Not applied in black-and-white landscape mode, which already uses \LARGE.
\ifwhiteblack\else
  \setminted{fontsize=\small, breaklines=true}
\fi

% unified tcolorboxes for articles
\ifwhiteblack
\tcbset{colback=white, colframe=black, fonttitle=\bfseries, boxrule=1pt}
\newtcolorbox{boxDef}[1][]{title={{\emoji📘} تعريف\ifx\\#1\\\else ~#1\fi :}}
\newtcolorbox{boxExercise}[1][]{title={{\emoji🧩} تمرين\ifx\\#1\\\else ~#1\fi :}}
\newtcolorbox{boxExample}[1][]{title={{\emoji📝} مثال\ifx\\#1\\\else ~#1\fi :}}
\newtcolorbox{boxNote}[1][]{title={{\emoji✨} ملاحظة\ifx\\#1\\\else ~#1\fi :}}
\newtcolorbox{boxAttention}[1][]{title={{\emoji🔔} تنبيه\ifx\\#1\\\else ~#1\fi :}}
\newtcolorbox{boxWarning}[1][]{title={{\emoji⚡} ملاحظة هامة\ifx\\#1\\\else ~#1\fi :}}
\newtcolorbox{boxSolution}[1][]{title={{\emoji✅} حل\ifx\\#1\\\else ~#1\fi :}}
\newtcolorbox{boxSymbol}[1][]{title={{\emoji🔣} رمز\ifx\\#1\\\else ~#1\fi :}}
\newtcolorbox{boxHint}[1][]{title={{\emoji💡} تلميح\ifx\\#1\\\else ~#1\fi :}}
\else
\tcbset{colback=white, colframe=black, fonttitle=\bfseries, boxrule=0.8pt}
\newtcolorbox{boxDef}[1][]{colback=blue!5!white,colframe=blue!75!black,
  title={{\emoji📘} تعريف\ifx\\#1\\\else ~#1\fi :}}
\newtcolorbox{boxExercise}[1][]{colback=cyan!5!white,colframe=cyan!70!black,
  title={{\emoji🧩} تمرين\ifx\\#1\\\else ~#1\fi :}}
\newtcolorbox{boxExample}[1][]{colback=yellow!5!white,colframe=orange!90!black,
  title={{\emoji📝} مثال\ifx\\#1\\\else ~#1\fi :}}
\newtcolorbox{boxNote}[1][]{colback=gray!10!white,colframe=black,
  title={{\emoji✨} ملاحظة\ifx\\#1\\\else ~#1\fi :}}
\newtcolorbox{boxAttention}[1][]{colback=magenta!10!white,colframe=magenta!80!black,
  title={{\emoji🔔} تنبيه\ifx\\#1\\\else ~#1\fi :}}
\newtcolorbox{boxWarning}[1][]{colback=red!5!white,colframe=red!75!black,
  title={{\emoji⚡} ملاحظة هامة\ifx\\#1\\\else ~#1\fi :}}
\newtcolorbox{boxSolution}[1][]{colback=green!5!white,colframe=green!60!black,
  title={{\emoji✅} حل\ifx\\#1\\\else ~#1\fi :}}
\newtcolorbox{boxSymbol}[1][]{colback=purple!5!white,colframe=purple!70!black,
  title={{\emoji🔣} رمز\ifx\\#1\\\else ~#1\fi :}}
\newtcolorbox{boxHint}[1][]{colback=teal!5!white,colframe=teal!60!black,
  title={{\emoji💡} تلميح\ifx\\#1\\\else ~#1\fi :}}
\fi


\ifwhiteblack
\tcbset{simplecode/.style={ colback=white, colframe=black, boxrule=1pt, arc=2pt, left=4pt,right=4pt,top=4pt,bottom=4pt}}
\else
\tcbset{simplecode/.style={ colback=gray!5, colframe=black!50, boxrule=0.4pt, arc=2pt, left=4pt,right=4pt,top=4pt,bottom=4pt}}
\fi
% Compact code boxes: tighter line spacing (the preamble uses onehalfspacing) and small vertical gaps.
% Must come AFTER the simplecode definition above.
\tcbset{simplecode/.append style={fontupper=\setstretch{0.95}, top=4pt, bottom=4pt, before skip=6pt, after skip=6pt}}
\newenvironment{boxCode}{\begin{tcolorbox}[simplecode]}{\end{tcolorbox}}

% Code cell used inside tables (keeps code left-to-right inside RTL tables).
% Use inside a fixed-width p{...} column: \begin{pycell}\begin{minted}...\end{minted}\end{pycell}
\newtcolorbox{pycell}{frame empty,colback=white,boxsep=0pt,left=2pt,right=2pt,top=2pt,bottom=2pt,width=\linewidth}

\newcolumntype{C}[1]{>{\centering\arraybackslash}p{#1}}

\makeatletter
\DTMnewdatestyle{dotted}{%
  \renewcommand{\DTMdisplaydate}[4]{\two@digits{##3}.\two@digits{##2}.##1}%
  \renewcommand{\DTMDisplaydate}{\DTMdisplaydate}%
}
\makeatother
\DTMsetdatestyle{dotted}

% redefine spaces after titles
\makeatletter
\renewcommand{\@maketitle}{%
  \begin{center}
    {\huge \bfseries \@title \par}%
    \vskip 0.2em % space between title and author
    {\large \@author \par}%
    % \vskip 0.2em % space between author and date
    % {\normalsize \@date \par}%
    {\small \hfill \begin{english}\DTMtoday\end{english} \par}%
  \end{center}
}
\makeatother

\fancyhf{} % clear default
\fancypagestyle{plain}{
  \fancyhf{}
  \fancyhead[L]{مدرسة التسامح الشاملة}
  % \fancyhead[L]{\includegraphics[height=1cm]{../../../images/logoTasamoh.png}}
  \fancyhead[R]{الأستاذ محمود اغبارية}
  \fancyfoot[C]{\thepage}
}

\fancyhead[L]{مدرسة التسامح الشاملة}
\fancyhead[R]{الأستاذ محمود اغبارية}
\fancyfoot[C]{\thepage}

\setcounter{tocdepth}{3} % only section subsection and subsubsection in TOC

\newcommand{\importpdfpage}[2]{%
    \thispagestyle{empty}
    \newgeometry{left=0cm,right=0cm,top=0cm,bottom=0cm}
    \includegraphics[page=#2,width=0.95\paperwidth,height=0.95\paperheight,keepaspectratio]{#1}
    \restoregeometry
}

\newcommand{\insertFullImg}[1]{%
    \noindent
    \makebox[\textwidth][c]{
        \includegraphics[width=0.9\paperwidth,keepaspectratio]{#1}%
    }%
}%

\usepackage{longtable}
\providecommand{\LTcaptype}{table}

% Code cell used inside tables (keeps code left-to-right inside RTL tables)
\newtcolorbox{pycell}{frame empty,colback=white,boxsep=0pt,left=2pt,right=2pt,top=2pt,bottom=2pt,width=\linewidth}

\ifwithsols
\title{حل ورقة تمرن - حلقات \textenglish{for} و \textenglish{while}}
\else
\title{ورقة تمرن - حلقات \textenglish{for} و \textenglish{while}}
\fi

\begin{document}

\maketitle
\thispagestyle{fancy}

%==============================================================
\section*{القسم الأول: حلقات \textenglish{for}}

\subsection*{سؤال 1}
اكتبوا \textbf{المخرج} الناتج عن كل مقطع برنامج في الجدول التالي.

{\small
\renewcommand{\arraystretch}{1.3}
\begin{longtable}{|c|p{0.45\textwidth}|p{0.32\textwidth}|}
\hline
\multicolumn{1}{|c|}{\textbf{رقم البند}} & \multicolumn{1}{c|}{\textbf{الكود في لغة \textenglish{python}}} & \multicolumn{1}{c|}{\textbf{المخرج الناتج}} \\
\hline
\endhead
1 &
\begin{pycell}
\begin{english}
\begin{minted}[fontsize=\small]{python}
for item in range(3):
    print(item + 1, end=" ")
\end{minted}
\end{english}
\end{pycell} & \rule{0pt}{1.3cm} \\
\hline
2 &
\begin{pycell}
\begin{english}
\begin{minted}[fontsize=\small]{python}
for item in range(3, 5):
    print(item)
\end{minted}
\end{english}
\end{pycell} & \rule{0pt}{1.3cm} \\
\hline
3 &
\begin{pycell}
\begin{english}
\begin{minted}[fontsize=\small]{python}
for item in range(6, 1, -2):
    print(item, end=" ")
\end{minted}
\end{english}
\end{pycell} & \rule{0pt}{1.3cm} \\
\hline
4 &
\begin{pycell}
\begin{english}
\begin{minted}[fontsize=\small]{python}
for item in range(-5, 6):
    print(item + 6)
\end{minted}
\end{english}
\end{pycell} & \rule{0pt}{1.3cm} \\
\hline
5 &
\begin{pycell}
\begin{english}
\begin{minted}[fontsize=\small]{python}
for i in range(16, 3, -4):
    print(i)
\end{minted}
\end{english}
\end{pycell} & \rule{0pt}{1.3cm} \\
\hline
6 &
\begin{pycell}
\begin{english}
\begin{minted}[fontsize=\small]{python}
for item in range(6):
    if item % 2 == 1:
        print(item * item)
\end{minted}
\end{english}
\end{pycell} & \rule{0pt}{1.3cm} \\
\hline
7 &
\begin{pycell}
\begin{english}
\begin{minted}[fontsize=\small]{python}
for item in range(9, 0, -3):
    print(item)
\end{minted}
\end{english}
\end{pycell} & \rule{0pt}{1.3cm} \\
\hline
8 &
\begin{pycell}
\begin{english}
\begin{minted}[fontsize=\small]{python}
for x in range(7, 8):
    print(x)
\end{minted}
\end{english}
\end{pycell} & \rule{0pt}{1.3cm} \\
\hline
9 &
\begin{pycell}
\begin{english}
\begin{minted}[fontsize=\small]{python}
for x in range(100, 120):
    if x % 2 == 0:
        print(x)
\end{minted}
\end{english}
\end{pycell} & \rule{0pt}{1.3cm} \\
\hline
10 &
\begin{pycell}
\begin{english}
\begin{minted}[fontsize=\small]{python}
for x in range(100, 120, 2):
    print(x)
\end{minted}
\end{english}
\end{pycell} & \rule{0pt}{1.3cm} \\
\hline
11 &
\begin{pycell}
\begin{english}
\begin{minted}[fontsize=\small]{python}
for x in range(1, 10):
    print(x - 1)
\end{minted}
\end{english}
\end{pycell} & \rule{0pt}{1.3cm} \\
\hline
12 &
\begin{pycell}
\begin{english}
\begin{minted}[fontsize=\small]{python}
for x in range(4, 5):
    print(x + 3)
\end{minted}
\end{english}
\end{pycell} & \rule{0pt}{1.3cm} \\
\hline
13 &
\begin{pycell}
\begin{english}
\begin{minted}[fontsize=\small]{python}
for x in range(-5, 6):
    print(x + 6)
\end{minted}
\end{english}
\end{pycell} & \rule{0pt}{1.3cm} \\
\hline
\end{longtable}
}

\ifwithsols
\begin{boxSolution}
في جميع البنود تُطبع كل قيمة في \textbf{سطر مستقل}، ما عدا البندين 1 و3 حيث تُطبع جميع القيم في سطر واحد بسبب استعمال \textenglish{end=" "}.

\begin{english}
\renewcommand{\arraystretch}{1.2}
\begin{tabular}{rcl}
1 & $\rightarrow$ & \texttt{1 2 3} \\
2 & $\rightarrow$ & \texttt{3, 4} \\
3 & $\rightarrow$ & \texttt{6 4 2} \\
4 & $\rightarrow$ & \texttt{1, 2, 3, \ldots, 11} \\
5 & $\rightarrow$ & \texttt{16, 12, 8, 4} \\
6 & $\rightarrow$ & \texttt{1, 9, 25} \\
7 & $\rightarrow$ & \texttt{9, 6, 3} \\
8 & $\rightarrow$ & \texttt{7} \\
9 & $\rightarrow$ & \texttt{100, 102, 104, \ldots, 118} \\
10 & $\rightarrow$ & \texttt{100, 102, 104, \ldots, 118} \\
11 & $\rightarrow$ & \texttt{0, 1, 2, \ldots, 8} \\
12 & $\rightarrow$ & \texttt{7} \\
13 & $\rightarrow$ & \texttt{1, 2, 3, \ldots, 11} \\
\end{tabular}
\end{english}
\end{boxSolution}
\fi

\clearpage
\subsection*{سؤال 2}
طُلب من التلاميذ أن يكتبوا برنامجًا يستقبل 5 أزواج أعداد ويطبع \textbf{عدد الأزواج} التي لهما نفس باقي القسمة على 2 (أي عددين زوجيين أو عددين فرديين).

\textbf{مثال:}
\begin{itemize}
    \item زوج الأعداد الأول: \textenglish{4 \ 8} $\leftarrow$ نفس باقي القسمة على 2، العددين زوجيين.
    \item زوج الأعداد الثاني: \textenglish{7 \ 66} $\leftarrow$ باقي القسمة على 2 مختلف.
    \item زوج الأعداد الثالث: \textenglish{3 \ 5} $\leftarrow$ نفس باقي القسمة على 2، العددين فرديين.
    \item زوج الأعداد الرابع: \textenglish{1 \ 8} $\leftarrow$ باقي القسمة على 2 مختلف.
    \item زوج الأعداد الخامس: \textenglish{1 \ 5} $\leftarrow$ نفس باقي القسمة على 2، العددين فرديين.
\end{itemize}
المخرج الناتج من البرنامج هو \textenglish{3}، أي عدد الأزواج ذات نفس باقي القسمة على 2 هي \textenglish{3}.

أكملوا مقطع الكود التالي بحيث يقوم بالمطلوب.

\begin{boxCode}
\begin{english}
\begin{minted}{python}
count = ______________
for item in range(______________):
    num1 = int(input("Please enter first number: "))
    num2 = int(input("Please enter second number: "))
    if ______________:
        count = ______________
print(______________)
\end{minted}
\end{english}
\end{boxCode}

\textbf{مخزن الكلمات:}
\begin{center}
\begin{english}
\texttt{count} \quad \texttt{count + 1} \quad \texttt{num1\%2 == num2\%2} \quad \texttt{5} \quad \texttt{0}
\end{english}
\end{center}

\ifwithsols
\begin{boxSolution}
\begin{english}
\begin{minted}{python}
count = 0
for item in range(5):
    num1 = int(input("Please enter first number: "))
    num2 = int(input("Please enter second number: "))
    if num1 % 2 == num2 % 2:
        count = count + 1
print(count)
\end{minted}
\end{english}
\end{boxSolution}
\fi

\clearpage

%==============================================================
\section*{القسم الثاني: حلقات \textenglish{while}}

الحلقة \textbf{اللانهائية} هي الحلقة التي شرط التوقف فيها لا يتحقق أبدًا، ولذلك الأوامر التي في جسم الحلقة تتنفذ إلى ما لا نهاية. أمامكم مقاطع برنامج تحتوي على حلقة \textenglish{while}. لكل مقطع برنامج، حدّدوا إذا كانت الحلقة لا نهائية أم نهائية.
إذا كانت الحلقة لا نهائية، أشيروا بـ \textenglish{X} في الخانة الملائمة، وإذا كانت نهائية اكتبوا ماذا سيكون المخرج (أو عدد المرات التي ستُنفّذ فيها الحلقة).

\subsection*{سؤال 1}

{\small
\renewcommand{\arraystretch}{1.3}
\begin{longtable}{|c|p{0.34\textwidth}|p{0.17\textwidth}|p{0.2\textwidth}|}
\hline
\multicolumn{1}{|c|}{\textbf{البند}} & \multicolumn{1}{c|}{\textbf{مقطع البرنامج}} & \multicolumn{1}{c|}{\textbf{حلقة لا نهائية}} & \multicolumn{1}{c|}{\textbf{عدد المرات (أو المخرج)}} \\
\hline
\endhead
1 &
\begin{pycell}
\begin{english}
\begin{minted}[fontsize=\small]{python}
num = 5
while num > 0:
    print(num)
    num = num - 1
\end{minted}
\end{english}
\end{pycell} & \rule{0pt}{1.7cm} & \\
\hline
2 &
\begin{pycell}
\begin{english}
\begin{minted}[fontsize=\small]{python}
num = 0
while num < 10:
    print(num)
    num = num + 3
\end{minted}
\end{english}
\end{pycell} & \rule{0pt}{1.7cm} & \\
\hline
3 &
\begin{pycell}
\begin{english}
\begin{minted}[fontsize=\small]{python}
num = 0
while num > 0:
    print(num)
    num = num + 1
\end{minted}
\end{english}
\end{pycell} & \rule{0pt}{1.7cm} & \\
\hline
4 &
\begin{pycell}
\begin{english}
\begin{minted}[fontsize=\small]{python}
i = 5
while i == 5:
    print(i)
    i = i / 10
\end{minted}
\end{english}
\end{pycell} & \rule{0pt}{1.7cm} & \\
\hline
5 &
\begin{pycell}
\begin{english}
\begin{minted}[fontsize=\small]{python}
i = 5
while i == 5:
    print(i)
    i = i % 10
\end{minted}
\end{english}
\end{pycell} & \rule{0pt}{1.7cm} & \\
\hline
\end{longtable}
}

\ifwithsols
\begin{boxSolution}
\begin{enumerate}[label=\arabic*., itemsep=1em]
    \item حلقة \textbf{نهائية}: المخرج \textenglish{\texttt{5, 4, 3, 2, 1}} وتُنفَّذ الحلقة 5 مرات.
    \item حلقة \textbf{نهائية}: المخرج \textenglish{\texttt{0, 3, 6, 9}} وتُنفَّذ الحلقة 4 مرات.
    \item حلقة \textbf{نهائية}: الشرط غير محقق منذ البداية، لذلك لا يوجد مخرج وتُنفَّذ الحلقة 0 مرات.
    \item حلقة \textbf{نهائية}: المخرج \textenglish{\texttt{5}} وتُنفَّذ الحلقة مرة واحدة، لأنّ قيمة \textenglish{i} تصبح \textenglish{\texttt{0.5}} بعد القسمة فينتفي الشرط.
    \item حلقة \textbf{لا نهائية}: قيمة \textenglish{i} تبقى دائمًا 5 لأنّ \textenglish{\texttt{i \% 10 = 5}}، ولذلك يبقى الشرط محققًا دائمًا.
\end{enumerate}
\end{boxSolution}
\fi

\subsection*{سؤال 2}

{\small
\renewcommand{\arraystretch}{1.3}
\begin{longtable}{|c|p{0.34\textwidth}|p{0.17\textwidth}|p{0.2\textwidth}|}
\hline
\multicolumn{1}{|c|}{\textbf{البند}} & \multicolumn{1}{c|}{\textbf{مقطع البرنامج}} & \multicolumn{1}{c|}{\textbf{حلقة لا نهائية}} & \multicolumn{1}{c|}{\textbf{عدد المرات (أو المخرج)}} \\
\hline
\endhead
1 &
\begin{pycell}
\begin{english}
\begin{minted}[fontsize=\small]{python}
num = 5
while num > 0:
    print(num)
    num = num - 1
\end{minted}
\end{english}
\end{pycell} & \rule{0pt}{1.7cm} & \\
\hline
2 &
\begin{pycell}
\begin{english}
\begin{minted}[fontsize=\small]{python}
num = 1
while num < 10:
    print(num)
    num = num + 3
\end{minted}
\end{english}
\end{pycell} & \rule{0pt}{1.7cm} & \\
\hline
3 &
\begin{pycell}
\begin{english}
\begin{minted}[fontsize=\small]{python}
num = 5
while num > 0:
    print(num)
    num = num + 1
\end{minted}
\end{english}
\end{pycell} & \rule{0pt}{1.7cm} & \\
\hline
4 &
\begin{pycell}
\begin{english}
\begin{minted}[fontsize=\small]{python}
num = 3
while num != 10:
    print(num)
    num = num + 3
\end{minted}
\end{english}
\end{pycell} & \rule{0pt}{1.7cm} & \\
\hline
5 &
\begin{pycell}
\begin{english}
\begin{minted}[fontsize=\small]{python}
num = 2
while num != 10:
    print(num)
    num = num + 2
\end{minted}
\end{english}
\end{pycell} & \rule{0pt}{1.7cm} & \\
\hline
\end{longtable}
}

\ifwithsols
\begin{boxSolution}
\begin{enumerate}[label=\arabic*., itemsep=1em]
    \item حلقة \textbf{نهائية}: المخرج \textenglish{\texttt{5, 4, 3, 2, 1}} وتُنفَّذ الحلقة 5 مرات.
    \item حلقة \textbf{نهائية}: المخرج \textenglish{\texttt{1, 4, 7}} وتُنفَّذ الحلقة 3 مرات.
    \item حلقة \textbf{لا نهائية}: قيمة \textenglish{num} تزداد باستمرار وتبقى دائمًا أكبر من 0.
    \item حلقة \textbf{لا نهائية}: قيم \textenglish{num} هي \textenglish{\texttt{3, 6, 9, 12, \ldots}} فتتجاوز 10 دون أن تساويها أبدًا.
    \item حلقة \textbf{نهائية}: المخرج \textenglish{\texttt{2, 4, 6, 8}} وتُنفَّذ الحلقة 4 مرات.
\end{enumerate}
\end{boxSolution}
\fi

\clearpage
\subsection*{سؤال 3}
اكتبوا ماذا سيكون \textbf{المخرج} في كل بند.

{\small
\renewcommand{\arraystretch}{1.3}
\begin{longtable}{|c|p{0.45\textwidth}|p{0.3\textwidth}|}
\hline
\multicolumn{1}{|c|}{\textbf{البند}} & \multicolumn{1}{c|}{\textbf{مقطع البرنامج}} & \multicolumn{1}{c|}{\textbf{المخرج}} \\
\hline
\endhead
1 &
\begin{pycell}
\begin{english}
\begin{minted}[fontsize=\small]{python}
i = 12
while i < 15:
    i = i + 1
    print(i)
\end{minted}
\end{english}
\end{pycell} & \rule{0pt}{1.7cm} \\
\hline
2 &
\begin{pycell}
\begin{english}
\begin{minted}[fontsize=\small]{python}
i = 6
while i > 0:
    if i % 2 == 1:
        print(i)
    i = i - 1
\end{minted}
\end{english}
\end{pycell} & \rule{0pt}{1.7cm} \\
\hline
\end{longtable}
}

\ifwithsols
\begin{boxSolution}
كل قيمة تُطبع في \textbf{سطر مستقل}.

\begin{english}
\renewcommand{\arraystretch}{1.2}
\begin{tabular}{rcl}
1 & $\rightarrow$ & \texttt{13, 14, 15} \\
2 & $\rightarrow$ & \texttt{5, 3, 1} \\
\end{tabular}
\end{english}
\end{boxSolution}
\fi

\clearpage
\subsection*{سؤال 4}
اكتبوا ماذا سيكون \textbf{المخرج} في كل بند.

{\small
\renewcommand{\arraystretch}{1.3}
\begin{longtable}{|c|p{0.45\textwidth}|p{0.3\textwidth}|}
\hline
\multicolumn{1}{|c|}{\textbf{البند}} & \multicolumn{1}{c|}{\textbf{مقطع البرنامج}} & \multicolumn{1}{c|}{\textbf{المخرج}} \\
\hline
\endhead
1 &
\begin{pycell}
\begin{english}
\begin{minted}[fontsize=\small]{python}
i = 1
while i < 10:
    print(i)
    i = i + 2
\end{minted}
\end{english}
\end{pycell} & \rule{0pt}{1.7cm} \\
\hline
2 &
\begin{pycell}
\begin{english}
\begin{minted}[fontsize=\small]{python}
i = 1
while i <= 30:
    if i % 3 == 0:
        print(i)
    i = i + 1
\end{minted}
\end{english}
\end{pycell} & \rule{0pt}{1.7cm} \\
\hline
3 &
\begin{pycell}
\begin{english}
\begin{minted}[fontsize=\small]{python}
i = 1
while i < 4:
    print('wow')
    i += 1
\end{minted}
\end{english}
\end{pycell} & \rule{0pt}{1.7cm} \\
\hline
4 &
\begin{pycell}
\begin{english}
\begin{minted}[fontsize=\small]{python}
x = 5
i = 1
while i < 4:
    x -= 1
    i += 1
    print(i, x)
\end{minted}
\end{english}
\end{pycell} & \rule{0pt}{1.7cm} \\
\hline
5 &
\begin{pycell}
\begin{english}
\begin{minted}[fontsize=\small]{python}
x = 0
while x < 9:
    print(x)
    x += 1
\end{minted}
\end{english}
\end{pycell} & \rule{0pt}{1.7cm} \\
\hline
6 &
\begin{pycell}
\begin{english}
\begin{minted}[fontsize=\small]{python}
x = 1
while x < 12:
    print(x)
    x += 1
\end{minted}
\end{english}
\end{pycell} & \rule{0pt}{1.7cm} \\
\hline
\end{longtable}
}

\ifwithsols
\begin{boxSolution}
كل قيمة (أو زوج قيم في البند 4) تُطبع في \textbf{سطر مستقل}.

\begin{english}
\renewcommand{\arraystretch}{1.2}
\begin{tabular}{rcl}
1 & $\rightarrow$ & \texttt{1, 3, 5, 7, 9} \\
2 & $\rightarrow$ & \texttt{3, 6, 9, \ldots, 30} \\
3 & $\rightarrow$ & \texttt{wow, wow, wow} \\
4 & $\rightarrow$ & \texttt{2 4, 3 3, 4 2} \\
5 & $\rightarrow$ & \texttt{0, 1, 2, \ldots, 8} \\
6 & $\rightarrow$ & \texttt{1, 2, 3, \ldots, 11} \\
\end{tabular}
\end{english}
\end{boxSolution}
\fi

\clearpage
%==============================================================
\section*{القسم الثالث: \textenglish{for} مقابل \textenglish{while}}

صِلوا كل مقطع برنامج في العمود أ بمقطع البرنامج في العمود ب الذي \textbf{يطبع المخرج نفسه}.

{\small
\renewcommand{\arraystretch}{1.3}
\begin{longtable}{|c|p{0.36\textwidth}|p{0.36\textwidth}|c|}
\hline
\multicolumn{1}{|c|}{} & \multicolumn{1}{c|}{\textbf{العمود أ}} & \multicolumn{1}{c|}{\textbf{العمود ب}} & \multicolumn{1}{c|}{} \\
\hline
\endhead
1 &
\begin{pycell}
\begin{english}
\begin{minted}[fontsize=\small]{python}
i = 1
while i <= 10:
    if i % 3 == 0:
        print(i)
    i = i + 1
\end{minted}
\end{english}
\end{pycell} &
\begin{pycell}
\begin{english}
\begin{minted}[fontsize=\small]{python}
for i in range(-4, 5, 2):
    print(i + 5)
\end{minted}
\end{english}
\end{pycell} & \textenglish{A} \\
\hline
2 &
\begin{pycell}
\begin{english}
\begin{minted}[fontsize=\small]{python}
for i in range(16, 3, -4):
    print(i)
\end{minted}
\end{english}
\end{pycell} &
\begin{pycell}
\begin{english}
\begin{minted}[fontsize=\small]{python}
i = 1
while i <= 3:
    print(i * 3)
    i = i + 1
\end{minted}
\end{english}
\end{pycell} & \textenglish{B} \\
\hline
3 &
\begin{pycell}
\begin{english}
\begin{minted}[fontsize=\small]{python}
i = 1
while i < 10:
    print(i)
    i = i + 2
\end{minted}
\end{english}
\end{pycell} &
\begin{pycell}
\begin{english}
\begin{minted}[fontsize=\small]{python}
i = 8
while i > 0:
    print(i * 2)
    i = i - 2
\end{minted}
\end{english}
\end{pycell} & \textenglish{C} \\
\hline
4 &
\begin{pycell}
\begin{english}
\begin{minted}[fontsize=\small]{python}
for i in range(-5, 6):
    print(i + 6)
\end{minted}
\end{english}
\end{pycell} &
\begin{pycell}
\begin{english}
\begin{minted}[fontsize=\small]{python}
for i in range(1, 12):
    print(i)
\end{minted}
\end{english}
\end{pycell} & \textenglish{D} \\
\hline
\end{longtable}
}

\vspace{1.5cm}
\begin{center}
\renewcommand{\arraystretch}{1.3}
\begin{tabular}{|p{2cm}|p{3cm}|p{6cm}|}
\hline
رقم العمود أ &  حرف العمود ب & المخرج \\
\hline
1 &  &  \\
\hline
2 &  &  \\
\hline
3 &  &  \\
\hline
4 &  &  \\
\hline
\end{tabular}
\end{center}

\ifwithsols
\begin{boxSolution}
\begin{english}
\begin{tabular}{ll}
1 $\rightarrow$ B & (\texttt{3, 6, 9}) \\
2 $\rightarrow$ C & (\texttt{16, 12, 8, 4}) \\
3 $\rightarrow$ A & (\texttt{1, 3, 5, 7, 9}) \\
4 $\rightarrow$ D & (\texttt{1, 2, \ldots, 11}) \\
\end{tabular}
\end{english}
\end{boxSolution}
\fi

\clearpage

\section*{القسم الرابع: أكمل الكود}

\begin{enumerate}[itemsep=3em]


\item يوجد في مصنع الخرز ماكينة تستقبل خرزات بألوان مختلفة. أمامكم مقطع كود هدفه أن يعدّ \textbf{عدد الخرزات الحمراء}.

مثلًا، إذا كان المدخل (من اليسار لليمين): \textenglish{red red green red green red yellow yellow}، يكون المخرج: \textenglish{The number of red beads is: 4}.

أكملوا الناقص في البرنامج باستخدام مخزن الكلمات.

\begin{boxCode}
\begin{english}
\begin{minted}{python}
count_red = 0
ans = "yes"
while ans == "yes":
    ______ = input("Please enter a color: ")
    if color == ______:
        count_red = __________
    ans = input("Do you want to continue? yes/no ")
print("The number of red beads is: ", __________)
\end{minted}
\end{english}
\end{boxCode}

\textbf{مخزن الكلمات:}
\begin{center}
\begin{english}
\begin{boxCode}
\texttt{count\_red + 1} \quad , \quad \texttt{color} \quad , \quad \texttt{count\_red} \quad , \quad \texttt{"red"}
\end{boxCode}
\end{english}
\end{center}

\ifwithsols
\begin{boxSolution}
\begin{english}
\begin{minted}{python}
count_red = 0
ans = "yes"
while ans == "yes":
    color = input("Please enter a color: ")
    if color == "red":
        count_red = count_red + 1
    ans = input("Do you want to continue? yes/no ")
print("The number of red beads is: ", count_red)
\end{minted}
\end{english}
\end{boxSolution}
\fi

\clearpage

\item يريد معلم أن يحسب عدد الطلاب الذين نجحوا في الامتحان (علامة 55 أو أكثر) من بين 8 طلاب في الصف.

أكملوا الناقص في البرنامج باستخدام مخزن الكلمات ليستقبل البرنامج 8 علامات ويحسب عدد الناجحين ثم يطبع عددهم.

\begin{boxCode}
\begin{english}
\begin{minted}{python}
passed = 0
for i in range(______):
    grade = int(input("Enter grade: "))
    if ______ >= 55:
        passed = __________
print("Number of passing students: ", passed)
\end{minted}
\end{english}
\end{boxCode}

\textbf{مخزن الكلمات:}
\begin{center}
\begin{english}
\begin{boxCode}
\texttt{passed + 1} \quad , \quad \texttt{grade} \quad , \quad \texttt{8} \quad , \quad \texttt{i}
\end{boxCode}
\end{english}
\end{center}

\ifwithsols
\begin{boxSolution}
\begin{english}
\begin{minted}{python}
passed = 0
for i in range(8):
    grade = int(input("Enter grade: "))
    if grade >= 55:
        passed = passed + 1
print("Number of passing students: ", passed)
\end{minted}
\end{english}
\end{boxSolution}
\fi


\item آلة حاسبة في سوبر ماركت تجمع أسعار المنتجات التي يشتريها الزبون. تتوقف الآلة عن استقبال الأسعار وتطبع المجموع الكلي عندما يتم إدخال السعر 0.

أكملوا الناقص في البرنامج باستخدام مخزن الكلمات.

\begin{boxCode}
\begin{english}
\begin{minted}{python}
total = ______
price = int(input("Enter price (0 to stop): "))
while price != ______:
    total = total + ______
    price = int(input("Enter price (0 to stop): "))
print("Total to pay is: ", total)
\end{minted}
\end{english}
\end{boxCode}

\textbf{مخزن الكلمات:}
\begin{center}
\begin{english}
\begin{boxCode}
\texttt{price} \quad , \quad \texttt{0} \quad , \quad \texttt{total} \quad , \quad \texttt{0}
\end{boxCode}
\end{english}
\end{center}

\ifwithsols
\begin{boxSolution}
\begin{english}
\begin{minted}{python}
total = 0
price = int(input("Enter price (0 to stop): "))
while price != 0:
    total = total + price
    price = int(input("Enter price (0 to stop): "))
print("Total to pay is: ", total)
\end{minted}
\end{english}
\end{boxSolution}
\fi


\item برنامج لخزنة ذكية يطلب من المستخدم إدخال كلمة المرور السريّة (\textenglish{"1234"}). يستمر البرنامج بطلب كلمة المرور طالما أن المستخدم يُدخل كلمة مرور خاطئة. عند إدخال كلمة المرور الصحيحة، يطبع البرنامج عبارة \textenglish{"Safe Opened!"}.

أكملوا الناقص في البرنامج باستخدام مخزن الكلمات.

\begin{boxCode}
\begin{english}
\begin{minted}{python}
password = input("Enter password: ")
while password ______ "1234":
    print("Wrong password, try again.")
    ______ = input("Enter password: ")
print(__________)
\end{minted}
\end{english}
\end{boxCode}

\textbf{مخزن الكلمات:}
\begin{center}
\begin{english}
\begin{boxCode}
\texttt{"Safe Opened!"} \quad , \quad \texttt{password} \quad , \quad \texttt{==} \quad , \quad \texttt{!=}
\end{boxCode}
\end{english}
\end{center}

\ifwithsols
\begin{boxSolution}
\begin{english}
\begin{minted}{python}
password = input("Enter password: ")
while password != "1234":
    print("Wrong password, try again.")
    password = input("Enter password: ")
print("Safe Opened!")
\end{minted}
\end{english}
\end{boxSolution}
\fi

\end{enumerate}

\vspace{1cm}
\begin{flushleft}
أرجو لكم وقتًا ممتعًا.

الأستاذ محمود اغبارية.
\end{flushleft}

\end{document}
